\scanpage{232}\mspage{87}
Je garderai les notations $\unsure{\sharp|\natural},\flat$ \unsure{du|des} \unsure{par.} \unsure{précédent|précédents}.
Quels sont les deux couples extrêmes ?

Partant de $A_0=\operatorname{Is}(M)$ (plus petit des éléments
de $P$), on trouve $\mathbb L_W(A_0)=\operatorname{Fl}(M)$ (plus grand
élément de $Q$), et $A_0^\sharp=\mathbb A_W(\operatorname{Fl}(M))=W_{\mathrm{univ.}}$,
ensemble des flèches qui sont dans $W$ et le
restent par tt chgt de base :
\[
\deleted{\operatorname{Is}(M)^\sharp=W_{\mathrm{univ}},\quad\operatorname{Is}(M)}
\]
\[
A_0^\sharp=W_{\mathrm{univ}},\quad
\deleted{\unsure{A|\alpha}}\beta_0(A_0)\overset{\text{déf}}{=}\mathbb L(A_0)=\operatorname{Fl}(M),
\tag*{(6.9.)}
\]
\margin{où $A_0=\operatorname{Is}(M)$, plus petit él. de $P$}

Inversement, partons de $L_0=\operatorname{Is}(M)$, le plus petit
élément de $Q$. On a $\mathbb A_W(L_0)=W$ ($W$
étant stable par composition et \unsure{contnt} les iso),
et $L_0^\flat=\mathbb L(W)=\operatorname{Fib}_W$, les $W$-fibrations.
\[
L_0^\flat=\operatorname{Fib}_W,\quad \alpha_0(L_0)=W,
\tag*{(6.10)}
\]
\margin{où $L_0=\operatorname{Is}(M)$, plus petit él. de $Q$.}

Donc les parties $\mathcal R$-closes à g. de $\operatorname{Fl} M$
se promènent entre $W_{\mathrm{univ}}$ et $W$, les
parties $\mathcal R$-closes à droite entre $\operatorname{Fib}_W$ et $\operatorname{Fl} M$.

\uline{Proposition} 6.1. Soient $A\in P$, $L\in Q$ tels
\llap{\smash{\rule[-4.9\baselineskip]{0.35pt}{5.2\baselineskip}}\hspace{0.5em}}que $A=\deleted{\unsure{\beta}}\alpha_0(L)$, $L=\beta_0(A)$ (i.e. $A=\mathbb A_W(L)$,
$L=\mathbb L_W(A)$). Pour que $L$ soit stable par
composition, il f. et s. que $A$ soit
stable par chgt. de base \deleted{pour} \unsure{par les|les}
flèches $f\in L$. (Plus généralement si $L\inserted{\in Q}$ est
un élément quelconque de $Q$, et si on suppose de
% archive sheet number: 232
